\uuid{pIc0}
\titre{Module/argument}
\auteur{bodin ; hanani ; mzari}
\organisation{exo7}
\licence{CC BY-NC-SA 4.0}
\datecreate{2020-01-20}
\niveau{L1}
\module{Algèbre}
\chapitre{Nombres complexes}
\sousChapitre{Forme cartésienne, forme polaire}
\difficulte{2}
\theme{Écritures algébrique et géométrique, module, argument}
\exo7qcmid{lille-1.nombres-complexes.008}

\contenu{
\begin{question}
\qclassification{104.01}
\qtags{motcle=module/argument}
Soit $z=\frac{\cos \theta + i \sin \theta}{\cos \phi - i \sin \phi}$, $\theta, \phi \in \Rr$.
Quelles sont les assertions vraies ?
\begin{answers}
    \bad{$|z|=2$}

    \good{$\arg z = \theta + \phi \, [2\pi]$}

    \good{ $z = \cos (\theta+\phi) + i \sin (\theta + \phi)$}

    \good{$|z|=1$}

\end{answers}
\begin{explanations}
Utiliser l'écrire trigonométrique et  la formule : $\frac{e^{i\theta}}{e^{-i\phi}}= e^{i(\theta + \phi)} $.
\end{explanations}
\end{question}
}
